\section{结论}
\label{sec:summary}

从 QCD 中规范不变算符乘积的欧氏算符乘积展开出发，我们推导了准 PDF、\spcorr 与赝 PDF 的因子化公式。这三个欧氏分布函数是相互关联的可观测量，都源自同一基本因子化。对\spcorr 而言，这一推导意味着式~(\ref{eq:ratio})中的比值确实按 $z^2$ 标度，但需要式~(\ref{eq:io-Q-fact})的小 $z^2$ 因子化公式来提取 PDF。
%
我们的因子化公式推导适用于重正化\spcorr 在任何方案下的定义。
这里使用的 OPE 还可以用来系统地推导式~(\ref{eq:momfact})幂次修正的因子化公式，那将涉及向高扭度部分子分布的匹配。（这些修正的数值相关性在文献~\cite{Chen:2016utp}中被考察。）
%
注意 LaMET 并不等价于来自 OPE 的展开：前者更为普遍，可用于其他量的格点计算，例如不存在简单 OPE 的 TMD-PDF。

我们对准 PDF 因子化公式的推导还验证了式~(\ref{eq:momfact})匹配系数中的部分子动量必须取 $p^z=yP^z$；与 $p^z=P^z$ 相比，这对 $\overline{\rm MS}$ 匹配结果造成相当大的差别（见图~\ref{fig:quasi}）。因此，在 LaMET 方法中计算 PDF 的格点计算应使用正确的 $p^z$。

作为对各分布之间关系及因子化公式的一个非平凡检验，我们考虑了 $\overline{\rm MS}$ 方案下的单圈结果。我们导出了修正后的 $\overline{\rm MS}$ 方案系数 $C$，见式~(\ref{eq:quasi-matching})，它使卷积积分收敛。我们还计算了出现在\spcorr 与赝 PDF 因子化中的 Wilson 系数 ${\cal C}$ 的单圈 $\overline{\rm MS}$ 结果，并提供了这些单圈修正的数值分析。单圈匹配系数 ${\cal C}$ 对赝 PDF 的效应比 $C$ 对准 PDF 的效应小，比较图~\ref{fig:quasi}与图~\ref{fig:pseudo}即可看出。鉴于处理格点数据的系统不确定性，利用相同的空间关联函数格点数据、同时采用准 PDF 与赝 PDF 两种方法来提取部分子分布函数是有潜在价值的。

存在若干种实施因子化公式、从\spcorr $\tilde Q$ 的格点数据计算 PDF 的方式，我们已在\sec{lattice}中讨论。人们总是必须使用短距离关联与大核子动量来减小高扭度修正。这限制了格点计算中有用数据点的数目，如\sec{lattice}所述。要在不作模型假设的情况下实现精确计算，强烈期望走向更细的格距以增加有效数据点的数目。
